% AulaTeX - maqueta inicial de construcción descendente
% Contrato futuro: el Agente puede investigar, redactar, evaluar y compilar a partir de este archivo.
\documentclass[12pt]{article}
\usepackage[utf8]{inputenc}
\usepackage[T1]{fontenc}
\usepackage[spanish]{babel}
\usepackage{enumitem}
\usepackage{hyperref}
\usepackage{longtable}
\usepackage{booktabs}
\usepackage{xcolor}

\definecolor{uanlBlue}{HTML}{00376C}
\definecolor{uanlGold}{HTML}{D4AF37}
\definecolor{uanlGreen}{HTML}{1B5E20}

\title{\textbf{Nodo Institucional: Universidad Autónoma de Nuevo León (UANL)}}
\author{AulaTeX Académico}
\date{\today}

\begin{document}
\maketitle

\noindent\textbf{Nodo editorial:} biologia-celular-fagro\\
\textbf{Padre editorial:} interinstitucional\\
\textbf{Modo de generación:} reforzar\\
\textbf{Destino:} UANL\\
\textbf{Contrato futuro:} ready-for-agent

\section*{Ingesta base y estado del nodo}
\begin{itemize}[leftmargin=*]
  \item Institución: Universidad Autónoma de Nuevo León (UANL).
  \item Dependencias primarias: Facultad de Agronomía (FAGRO), Facultad de Ingeniería Mecánica y Eléctrica (FIME).
  \item Plataforma de gestión curricular: Nexus UANL / AulaTeX.
  \item Plantillas maestras activas: \texttt{reporte-uanl.tex}, \texttt{presentacion-uanl.tex}, \texttt{plantilla-actividad-uanl.tex}.
  \item Repositorio bibliográfico institucional: \texttt{bibliografia-uanl.bib}.
\end{itemize}

\section*{Objetivo formativo}
\begin{itemize}[leftmargin=*]
  \item Estructurar una base editorial e institucional robusta para la producción de evidencias de aprendizaje, reportes de laboratorio, proyectos integradores (PIA) y defensas orales en la UANL.
  \item Normalizar los estándares visuales (paleta Azul UANL, Oro y Verde Agronomía) y tipográficos (Arial/Helvetica, 1.5 de interlineado).
  \item Facilitar la generación automatizada y trazable de documentos científicos con aparato crítico en APA 7ma edición.
\end{itemize}

\section*{Competencias del Modelo Educativo UANL}
\begin{itemize}[leftmargin=*]
  \item \textbf{Formación integral:} Capacidad para interpretar fenómenos científicos complejos y resolver problemas del entorno socioproductivo.
  \item \textbf{Investigación y análisis crítico:} Trazabilidad metodológica, uso de fuentes indexadas primarias y contraste riguroso de hipótesis.
  \item \textbf{Comunicación profesional y científica:} Producción de documentos y defensas orales síncronas con estructura clara, solvencia léxica y soporte visual vectorial.
\end{itemize}

\section*{Arquitectura de entregables normalizados}
\begin{center}
\begin{tabular}{p{0.25\linewidth} p{0.35\linewidth} p{0.30\linewidth}}
\toprule
\textbf{Entregable} & \textbf{Plantilla Fuente} & \textbf{Salida Verificada} \\
\midrule
Reporte / PIA Escrito & \texttt{reporte-uanl.tex} & \texttt{reporte-uanl.pdf} \\
Defensa Oral Beamer & \texttt{presentacion-uanl.tex} & \texttt{presentacion-uanl.pdf} \\
Actividad Formativa & \texttt{plantilla-actividad-uanl.tex} & Formato XeLaTeX / PDF \\
Maqueta del Nodo & \texttt{maqueta.tex} & \texttt{maqueta.pdf} \\
\bottomrule
\end{tabular}
\end{center}

\section*{Criterios de evaluación y calidad}
\begin{itemize}[leftmargin=*]
  \item \textbf{Cero errores de compilación:} Procesamiento limpio mediante \texttt{latexmk} / \texttt{XeLaTeX}.
  \item \textbf{Aparato crítico estricto:} Integración de citas autor-año y bibliografía completa sin fuentes huérfanas.
  \item \textbf{Transferencia disciplinar:} Cada entregable debe incluir un apartado explícito de aplicación práctica en la profesión.
\end{itemize}

\end{document}
